\documentclass[12pt,a4paper,twoside]{report}

\usepackage[T1]{fontenc}
\usepackage[utf8]{inputenc}
\usepackage{tgtermes}
\usepackage[a4paper,top=2.1cm,bottom=2.1cm,left=2.3cm,right=2.3cm]{geometry}
\usepackage{setspace}
\setstretch{1.12}
\usepackage{ragged2e}
\justifying
\setlength{\parindent}{0.5cm}
\setlength{\parskip}{0pt}
\setlength{\emergencystretch}{3em}
\raggedbottom
\renewcommand{\topfraction}{0.92}
\renewcommand{\bottomfraction}{0.9}
\renewcommand{\textfraction}{0.06}
\renewcommand{\floatpagefraction}{0.72}
\setcounter{topnumber}{3}
\setcounter{bottomnumber}{3}
\setcounter{totalnumber}{5}

\usepackage{fancyhdr}
\usepackage{lastpage}
\pagestyle{fancy}
\fancyhf{}
\fancyfoot[R]{\thepage\ /\ \pageref{LastPage}}
\renewcommand{\headrulewidth}{0pt}

\usepackage{titlesec}
\titleformat{\chapter}[display]{\bfseries\Large}{\chaptername\ \thechapter}{0.4cm}{\Huge}
\titleformat{\section}{\bfseries\Large}{\thesection}{0.2cm}{}
\titleformat{\subsection}{\bfseries\large}{\thesubsection}{0.2cm}{}
\titleformat{\subsubsection}{\bfseries\normalsize}{\thesubsubsection}{0.2cm}{}

\usepackage{caption}
\captionsetup{font=small,labelfont=bf}
\usepackage{graphicx}
\usepackage{float}
\usepackage{booktabs}
\usepackage{longtable}
\usepackage{tabularx}
\usepackage{array}
\usepackage{xcolor}
\usepackage{enumitem}
\setlist[itemize]{noitemsep,topsep=0.25em}
\setlist[enumerate]{noitemsep,topsep=0.25em}
\usepackage{tikz}
\usetikzlibrary{arrows.meta,positioning,shapes.geometric,shapes.misc,fit,calc}
\PassOptionsToPackage{obeyspaces}{url}
\usepackage{xurl}
\usepackage{adjustbox}

% --- Code, configuration, and command-output listings ---
\usepackage{listings}
\definecolor{codebg}{RGB}{248,249,250}
\definecolor{codeframe}{RGB}{206,212,218}
\definecolor{codekw}{RGB}{37,99,235}
\definecolor{codestr}{RGB}{22,128,61}
\definecolor{codecmt}{RGB}{107,114,128}
\definecolor{codeln}{RGB}{150,156,164}
\definecolor{termbg}{RGB}{245,246,248}
\lstdefinestyle{aria}{%
  backgroundcolor=\color{codebg},
  basicstyle=\ttfamily\footnotesize,
  keywordstyle=\color{codekw}\bfseries,
  stringstyle=\color{codestr},
  commentstyle=\color{codecmt}\itshape,
  numberstyle=\tiny\color{codeln},
  numbers=left, numbersep=8pt,
  frame=single, rulecolor=\color{codeframe},
  breaklines=true, breakatwhitespace=false,
  showstringspaces=false, tabsize=2,
  captionpos=b, xleftmargin=14pt, xrightmargin=4pt,
  aboveskip=0.4em, belowskip=0.25em,
  upquote=true, columns=fullflexible, keepspaces=true,
}
\lstdefinestyle{ariaterm}{%
  style=aria, backgroundcolor=\color{termbg},
  numbers=none, keywordstyle=\color{codekw},
}
\renewcommand{\lstlistingname}{Listing}
% YAML language definition for playbook/config excerpts
\lstdefinelanguage{yaml}{%
  keywords={true,false,null,yes,no,on,off},
  keywordstyle=\color{codekw}\bfseries,
  comment=[l]{\#}, commentstyle=\color{codecmt}\itshape,
  morestring=[b]', morestring=[b]", stringstyle=\color{codestr},
  literate={:}{{\textcolor{codekw}{:}}}1,
  sensitive=true,
}
% INI / dotenv-style configuration definition
\lstdefinelanguage{conf}{%
  comment=[l]{\#}, commentstyle=\color{codecmt}\itshape,
  morestring=[b]", stringstyle=\color{codestr},
  morekeywords={true,false,enabled,disabled},
  keywordstyle=\color{codekw},
  sensitive=false,
}

\usepackage{hyperref}
\hypersetup{
  colorlinks=true,
  linkcolor=black,
  citecolor=black,
  urlcolor=blue
}
\usepackage[backend=biber,style=ieee]{biblatex}
\addbibresource{bibliography.bib}

\newcolumntype{Y}{>{\raggedright\arraybackslash}X}
\newcommand{\todo}[1]{\textbf{[TODO: #1]}}
\newcommand{\project}{ARIA}
\newcommand{\projectlong}{Adaptive Response Intelligence Automation}
% url-style commands: robust (safe in captions / List of Figures) and
% breakable (long paths wrap instead of overflowing). Arguments use raw
% underscores. Defined with the url package so special characters and line
% breaks are handled like the original \path.
\urlstyle{tt}
\DeclareRobustCommand{\codepath}[1]{\nolinkurl{#1}}
\DeclareRobustCommand{\sourcefile}[1]{\nolinkurl{#1}}

\newcommand{\visualplaceholder}[3]{%
\begin{figure}[H]
\centering
\fbox{%
\begin{minipage}[c][#2][c]{0.88\textwidth}
\centering
\textbf{#1}\\[0.25cm]
\small #3
\end{minipage}}
\caption{#1}
\end{figure}
}

\newcommand{\screenshotplaceholder}[2]{%
\begin{figure}[H]
\centering
\fbox{%
\begin{minipage}[c][3.2cm][c]{0.88\textwidth}
\centering
\textbf{#1}\\[0.2cm]
\small #2
\end{minipage}}
\caption{#1}
\end{figure}
}


\newcommand{\reservedblock}[3]{%
\par\vspace{0.2cm}
\noindent\fbox{%
\begin{minipage}[c][#2][c]{0.94\textwidth}
\centering
\textbf{#1}\\[0.15cm]
\small #3
\end{minipage}}
\par\vspace{0.3cm}
}

\newcommand{\diagramimage}[2]{%
\begin{figure}[htbp]
\centering
\includegraphics[width=\textwidth,height=0.72\textheight,keepaspectratio]{assets/diagrams/#1}
\caption{#2}
\end{figure}
}
\newcommand{\diagramimagecompact}[2]{%
\begin{figure}[htbp]
\centering
\includegraphics[width=\textwidth,height=0.42\textheight,keepaspectratio]{assets/diagrams/#1}
\caption{#2}
\end{figure}
}
% Labeled diagram figure for UML/Mermaid renders (height-capped, cross-referenceable)
\newcommand{\umlfigure}[3]{%
\begin{figure}[htbp]
\centering
\includegraphics[width=\textwidth,height=0.50\textheight,keepaspectratio]{assets/diagrams/#1}
\caption{#2}\label{#3}
\end{figure}
}
\newcommand{\umlfigurecompact}[3]{%
\begin{figure}[htbp]
\centering
\includegraphics[width=\textwidth,height=0.46\textheight,keepaspectratio]{assets/diagrams/#1}
\caption{#2}\label{#3}
\end{figure}
}
\newcommand{\cloudshot}[3]{%
\begin{figure}[H]
\centering
\fbox{\includegraphics[width=#3,height=0.40\textheight,keepaspectratio]{assets/cloud/#1}}
\caption{#2}
\end{figure}
}
% --- Real product screenshots ---
% Main-chapter screenshot: framed, cross-referenceable, height-capped.
\newcommand{\screenshot}[3]{%
\begin{figure}[H]
\centering
\fbox{\includegraphics[width=0.92\textwidth,height=0.55\textheight,keepaspectratio]{assets/screenshots/#1}}
\caption{#2}\label{#3}
\end{figure}
}
% Compact screenshot for appendix evidence gallery (two fit per page region).
\newcommand{\screenshotc}[3]{%
\begin{figure}[H]
\centering
\fbox{\includegraphics[width=0.82\textwidth,height=0.40\textheight,keepaspectratio]{assets/screenshots/#1}}
\caption{#2}\label{#3}
\end{figure}
}
% Two screenshots side by side, one shared caption / LoF line.
\newcommand{\pairfig}[4]{%
\begin{figure}[H]
\centering
\fbox{\includegraphics[width=0.48\textwidth,height=0.33\textheight,keepaspectratio]{assets/screenshots/#1}}\hfill
\fbox{\includegraphics[width=0.48\textwidth,height=0.33\textheight,keepaspectratio]{assets/screenshots/#2}}
\caption{#3}\label{#4}
\end{figure}
}
% Four screenshots (2x2), one shared caption / LoF line, for dense tab evidence.
\newcommand{\quadfig}[6]{%
\begin{figure}[H]
\centering
\fbox{\includegraphics[width=0.47\textwidth,height=0.29\textheight,keepaspectratio]{assets/screenshots/#1}}\hfill
\fbox{\includegraphics[width=0.47\textwidth,height=0.29\textheight,keepaspectratio]{assets/screenshots/#2}}\\[2pt]
\fbox{\includegraphics[width=0.47\textwidth,height=0.29\textheight,keepaspectratio]{assets/screenshots/#3}}\hfill
\fbox{\includegraphics[width=0.47\textwidth,height=0.29\textheight,keepaspectratio]{assets/screenshots/#4}}
\caption{#5}\label{#6}
\end{figure}
}

